	
	
	% Source: AI-generated/holomorphic_invariants_no43_vs_engel.tex.
	% Its X is our Y; its fibres q,r,c are our fibres 1,0,infinity.
	The following is a working summary of the calculations in the supplied
	notes, with sketches of the arguments, not a new verification. We use
	the analytic models and topological conclusions of the preceding
	sections. Put $G=(Y_0)_{\mathrm{red}}$, so that $Y_0=4G$.
	
	\paragraph{Canonical bundle and pluricanonical systems.}
	The calculation in the notes gives
	\[
	K_Y\simeq f^*\ko_{\IP^1}(-1)\simeq\ko_Y(-4G).
	\]
	On the open contracting quotient, wedge a local two-form on $S$
	with the logarithmic differential in the line-bundle direction.
	This is invariant under the contracting action. At the
	$\III^*$-fibre the pole of order three of the relative elliptic
	differential on the order-four cover cancels the zero of the base
	differential. The twisted deck action preserves the resulting
	three-form, which descends through the free quotient. At the
	retained cusp it extends as on the original surface, and at the
	Mumford fibre the toric calculation gives a nowhere vanishing
	three-form, also along the double and triple loci. Thus there is
	no further correction to the canonical bundle.
	
	Since $f_*\ko_Y=\ko_{\IP^1}$, the projection formula gives, for $m\geq1$,
	\[
	h^0(Y,K_Y^{\otimes m})=0,\qquad
	h^0(Y,K_Y^{\otimes(-m)})=m+1.
	\]
	In particular $\kappa(Y)=-\infty$, whereas $\kappa(-K_Y)=1$.
	The anticanonical system is a base-point-free pencil and recovers
	$f$, up to an automorphism of the base.
	
	For comparison, write $g\colon Y'\to\IP^1$ for Engel's family and
	$B_4$ for its reduced multiplicity-four fibre. The corresponding
	calculation gives
	\[
	K_{Y'}\simeq\ko_{Y'}(-2B_4),\qquad
	K_{Y'}^{\otimes2}\simeq g^*\ko_{\IP^1}(-1),\qquad
	h^0(Y',K_{Y'}^{\otimes(-m)})=\lfloor m/2\rfloor+1,
	\]
	see \cite[Proposition~9.11]{Alpoge2026}. Thus $|-K_{Y'}|$ consists
	of $2B_4$ and $|-2K_{Y'}|$ recovers $g$. In particular
	$h^0(Y,-K_Y)=2$, whereas $h^0(Y',-K_{Y'})=1$. This distinguishes
	the two families up to biholomorphism without using algebraic
	reduction or the still unknown Hodge numbers. It does not by itself
	separate their deformation components.
	
	\paragraph{Cohomology of the structure sheaf.}
	The recorded direct images are
	\[
	f_*\ko_Y=\ko_{\IP^1},\qquad
	R^1f_*\ko_Y\simeq\ko_{\IP^1}\oplus\ko_{\IP^1}(-1),\qquad
	R^2f_*\ko_Y\simeq\ko_{\IP^1}(-1).
	\]
	Here one must use the scheme-theoretic fibres, including $4G$.
	Adjunction gives trivial dualising sheaf on each fibre, and their
	structure-sheaf cohomology has dimensions $(1,2,1)$. For $4G$ the
	degree-zero calculation is
	$(\mathbb C[s]/(s^4))^{\mu_4}=\mathbb C$; duality and constancy
	of the Euler characteristic give the other two dimensions.
	Grauert's theorem and relative duality then give local freeness
	and the formula for $R^2f_*\ko_Y$.
	
	The quotient character $\xi$ supplies a nowhere zero section of
	$R^1f_*\ko_Y$. Its nonvanishing is checked on the smooth torus
	cover at $0$, on the elliptic curve over the cusp at $1$, and by
	the normalization complex of the Mumford fibre at $\infty$.
	Riemann--Roch gives degree $-1$ for this rank-two bundle; the
	resulting extension of $\ko(-1)$ by $\ko$ splits because
	$H^1(\IP^1,\ko(1))=0$. Coherent Leray now gives
	\[
	(h^0(Y,\ko_Y),h^1(Y,\ko_Y),h^2(Y,\ko_Y),h^3(Y,\ko_Y))
	=(1,1,0,0).
	\]
	For the original family these computations are in
	\cite[Proposition~9.13]{Alpoge2026}.
	
